% SOURCE_AUTHORITY: SGA5_LNM589_SCAN, PDF pages 5--7 = printed Roman pages V--VII (scan-only authority).
% SOURCE_DOCUMENT_SHA256: B256EBD072A8C68209518412A263C9289C6F1854A346733D86F885930D5FE6CA
Después de los preliminares técnicos de la exposición V sobre los sistemas proyectivos ádicos, en la exposición VI se define la noción de $\mathbb Z_\ell$-haz constructible (resp., de haz constante torcido constructible, es decir, «liso» en la terminología de Deligne ([1], (SGA~$4^{1/2}$ Arcata))), y se extiende a los $\mathbb Z_\ell$-haces constructibles, por paso al límite a partir de los $\mathbb Z/\ell^n$-haces, el formalismo de las imágenes directas superiores con soportes propios. Sin embargo, la exposición VI no da una definición de una categoría derivada de los $\mathbb Z_\ell$-haces, provista de las operaciones $\lten$ y $\uRHom$ que prolonguen aquellas de las que se dispone en la categoría derivada de los $\mathbb Z/\ell^n$-haces. \emph{A fortiori}, tampoco se considera la extensión a los $\mathbb Z_\ell$-haces del formalismo de dualidad de SGA~4 XVIII. Esta cuestión era el objeto de la tesis (no publicada) de Jouanolou. El lector podrá consultar, no obstante, el comienzo de [1], donde se definen los $\uExt^i$ de los $\mathbb Z_\ell$-haces.

La exposición VII, independiente del resto del seminario, es sin duda una de las más interesantes. Contiene a) el cálculo de la cohomología de algunas variedades clásicas (espacios afines, espacios proyectivos y variedades de banderas); b) una exposición de la teoría de las clases de Chern en cohomología étale; c) una demostración de la «fórmula de autointersección» $i^*i_*x=x\,c_d(N)$ en cohomología étale (resp., en el anillo de Chow, demostración debida en este caso a Mumford), así como diversas aplicaciones de esta fórmula (cálculo de la cohomología de las variedades obtenidas por explosión y fórmula de Gauss--Bonnet).

% SOURCE_EMENDATION: printed Roman p. V repeats «L'exposé VII» here. The subject (finiteness conditions and Grothendieck groups), the printed contents, and the dependency diagram identify this paragraph as Exposé VIII; Spanish corrects the numeral accordingly.
La exposición VII contiene sorites sobre las condiciones de finitud en las categorías derivadas y los grupos de Grothendieck. Estas cuestiones se retoman en un marco más general en (SGA~6 I, IV).

La exposición IX, de J.-P. Serre, titulada «Introducción a la teoría de Brauer», fue publicada en [4] y no se reproduce en este volumen. Uno de sus principales resultados, debido a Swan y relativo a la existencia de un módulo proyectivo cuyo carácter es el «carácter de Swan», se utiliza en la exposición X, donde se demuestra una fórmula de Euler--Poincaré para un haz sobre una curva, ligeramente más general que la que figura en la exposición de Raynaud [3].

% SOURCE_PAGE_BOUNDARY: end of printed Roman p. V.
\clearpage

La exposición XI, redactada por I. Bucur, no existe, porque se perdió durante una mudanza y el autor no conservaba ninguna copia. Estaba dedicada a la teoría de Grothendieck de las trazas conmutativas, que generaliza la de Stallings [5]. Una versión más elaborada de esta teoría, en un marco de teoría de haces, se expone en (III~B 5, 6). En las exposiciones XII y III~B se dan aplicaciones a fórmulas de Lefschetz. La fórmula de las trazas de la exposición XII se demuestra con independencia de la fórmula general de la exposición III, pero en (III~B 6) se muestra que los términos locales que figuran en ella son precisamente los de la fórmula general, y que esta última la implica.

La ausencia de la exposición XIII se debe únicamente a una razón de numeración: inicialmente se había previsto que el contenido de las exposiciones X a XII se extendiera a cuatro exposiciones.

La exposición XIV (numerada a veces también XV) contiene generalidades sobre la correspondencia de Frobenius (entre ellas, la comparación entre el Frobenius «geométrico» y el Frobenius «aritmético») y, a partir de la fórmula de las trazas de la exposición XII (o III~B), da la demostración de Grothendieck de la racionalidad de las funciones $L$: los dos puntos importantes son la reducción al caso de las curvas y el paso al límite que permite reducirse a una fórmula para $\mathbb Z/\ell^n$-haces.

Durante el seminario oral, Grothendieck había pronunciado una exposición sobre los problemas abiertos y enunciado algunas conjeturas. Esta exposición, que debía clausurar el seminario, no fue desgraciadamente redactada, como tampoco lo fue su bellísima exposición introductoria, en la que pasaba revista a las fórmulas de Euler--Poincaré y de Lefschetz en diversos contextos (topológico, analítico complejo y algebraico).

Agradezco a P. Deligne que me convenciera de redactar, en una nueva versión de la exposición III, una demostración de la fórmula de Lefschetz--Verdier, eliminando así uno de los obstáculos para la publicación de este seminario. Le estoy muy agradecido por las mejoras de redacción que me sugirió,

% SOURCE_PAGE_BOUNDARY: the printed sentence continues from Roman p. VI to p. VII here.
\clearpage
\noindent y por la ayuda que me prestó en la preparación de este volumen para el editor. Agradezco asimismo a la Sra.~Lièvremont, que mecanografió con esmero, y en muy poco tiempo, las exposiciones III y III~B.

No puedo terminar esta introducción sin rendir homenaje a la memoria de I. Bucur, fallecido de cáncer en septiembre de 1976. Quienes, como yo, tuvimos el privilegio de tenerlo por amigo no olvidamos su exquisita gentileza, su humor risueño y el encanto de su conversación.

\begin{flushright}
París, 19 de febrero de 1977\\[1.2em]
Luc Illusie
\end{flushright}
